\section{Πλαίσιο και Αντικείμενο της Εργασίας}

Η ηλεκτρονική μάθηση (\ENG{E-Learning}) αποτελεί πλέον βασικό τμήμα των σύγχρονων εκπαιδευτικών και εταιρικών
περιβαλλόντων. Σημαντικό ρόλο σε αυτή την εξέλιξη διαδραματίζουν τα πρότυπα διαλειτουργικότητας εκπαιδευτικού
περιεχομένου, τα οποία επιτρέπουν τη μεταφερσιμότητα, την επαναχρησιμοποίηση και τη μακροχρόνια διατήρηση των
μαθημάτων.

Στο πλαίσιο αυτό, το πρότυπο \ENG{SCORM (Sharable Content Object Reference Model)} εξακολουθεί να αποτελεί την
πιο διαδεδομένη πρακτική λύση για την εκτέλεση και την παρακολούθηση μεγάλου όγκου υφιστάμενου εκπαιδευτικού
περιεχομένου. Παρότι έχουν εμφανιστεί νεότερα πρότυπα, μεγάλο μέρος των οργανισμών εξακολουθεί να βασίζεται σε
υλικό που έχει ήδη παραχθεί σύμφωνα με τις εκδόσεις \ENG{SCORM 1.2} και \ENG{SCORM 2004}.

Η παρούσα διπλωματική εργασία εξετάζει το πρόβλημα της εκτέλεσης \ENG{SCORM} από την οπτική της ανάπτυξης ενός
αυτόνομου \ENG{SCORM Engine}. Στόχος είναι η δημιουργία μιας πλατφόρμας που μπορεί να ενσωματωθεί σε διαφορετικά
πληροφοριακά συστήματα μάθησης, παρέχοντας έναν σαφή μηχανισμό εκτέλεσης μαθημάτων και καταγραφής προόδου.

Το κύριο παραδοτέο της εργασίας είναι το σύνολο του πηγαίου κώδικα, το οποίο φιλοξενείται σε
\ENG{GitHub Organization} και αποτελεί την πλήρη υλοποίηση της προτεινόμενης πλατφόρμας. Ο κώδικας συνοδεύεται από τα απαραίτητα
εργαλεία ανάπτυξης και τεκμηρίωσης ώστε το σύστημα να μπορεί να εγκατασταθεί και να εκτελεστεί ως ενιαίο
οικοσύστημα υπηρεσιών.

Παράλληλα με την υλοποίηση, πραγματοποιείται θεωρητική ανάλυση των προτύπων \ENG{E-Learning} και των τεχνολογικών
παραμέτρων που επηρεάζουν τον σχεδιασμό ενός σύγχρονου μηχανισμού εκτέλεσης \ENG{SCORM}.